\section{서론}
\label{sec:introduction}

% P1: 맥락, 관측 수단, 과제
이동하는 다수의 에이전트가 남기는 대규모 시공간 궤적은 사건 및 환경 기록과 결합될 때 에이전트 사이의 상호작용과 환경 변화까지 함께 분석할 수 있게 한다~\cite{Zheng15}.
\rev{배틀로얄 게임은 활동 영역이 줄어드는 환경의 생존 경쟁을 이러한 데이터로 대규모로 기록한다.}
PUBG: BATTLEGROUNDS(이하 PUBG)에서는 최대 100명의 플레이어가 1인, 2인 또는 4인 팀으로 경쟁하며, 안전구역은 정해진 단계에 따라 줄어들고 구역 밖의 플레이어는 지속적인 피해를 입는다.
공개 텔레메트리는 모든 플레이어의 위치와 교전, 아이템, 구역 사건을 담고 있어~\cite{PUBGAPI}, 같은 규칙에서 반복되는 경기로 줄어드는 영역 아래의 생존을 관찰할 수 있다.
\rev{새 안전구역이 공개되는 단계 경계마다 각 팀은 이동과 교전을 새로 정해야 한다.
본 연구는 그 경계에서 생존한 팀이 다음 단계가 끝나기 전에 탈락할 조건부 확률, 곧 다음 단계 탈락 위험을 분석하며, 이 양은 경기 복기, 중계 분석, 구역 규칙 검토~\cite{PUBGBlueZone2026}에 쓰일 수 있다.}
\rev{이는 모든 팀의 기록에 같은 절단 시각을 적용한 관전자 관점의 위험 예측으로, 새 구역을 본 직후 각 팀의 의사결정을 지원하는 것과는 구분된다.}

% P2: 난점과 기존 설계의 한계
\rev{그러나 이 위험은 줄어드는 영역의 기하와 주변 팀과의 상호작용에 함께 달려 있고 그 수준도 단계마다 다르며, 예측 시점 이후의 정보가 입력에 섞이면 결과가 누설된다.}
\rev{기존 PUBG 예측 연구는 이 문제를 서로 다른 단위와 대상으로 다룬다. Sagong 등은 단계별 활동 점수로 개별 플레이어의 생존을 분류하거나~\cite{Sagong25a} 전략 시퀀스로 승리를 예측하고~\cite{Sagong25b},} KRAFTON의 예측기~\cite{KRAFTON}는 팀 간 Transformer로 탈락 순서의 상대위험을 학습하여 우승확률로 변환한다.
그러나 상대위험 점수는 정해진 다음 단계 안의 탈락확률과 다른 추정 대상이며, 그 학습 목적만으로는 단계마다 공통적인 위험 수준이 식별되지 않는다.
또한 두 설계는 분석 단위, 데이터, 평가 지표가 모두 달라, 어떤 설계 선택이 성능을 좌우하는지 보고된 결과만으로는 알 수 없다.

% P3: 생태학적 관점
\rev{영역이 줄어드는 환경의 생존은} 서식지 손실 아래에서 집단의 존속을 다루는 생태학의 문제와 맞닿아 있다~\cite{Fahrig03}.
본 연구는 안전구역의 축소를 서식지의 단계적 감소로, 팀의 탈락을 집단의 소멸로 대응시켜, 다음 단계 탈락 위험을 단계를 시간 단위로 하는 이산시간 hazard로 정식화한다\rev{(\ref{sec:method:samples}절)}.
각 팀의 상태는 구역 경계에 대한 노출, \rev{이동생태학에서 착안한} 이동 부담~\cite{Nathan08}, 주변 경쟁 밀도, 남은 팀원 수로 구성한 서식지 맥락으로 표현한다.

% P4: 연구 질문과 접근
이 관점은 세 가지 질문을 제기한다.
\begin{itemize}
    \item \textbf{RQ1:}
    \rev{서식지 축소 관점의 네 요인(구역 경계 노출, 이동 부담, 주변 경쟁 밀도, 남은 팀원 수)은 다음 단계 탈락 위험과 어떻게 관련되며, 그 관계는 단계와 코호트에 걸쳐 일관되는가?}

    \item \textbf{RQ2:}
    \rev{단계 조건부 hazard 목적과 서식지 맥락을 결합한 설계는 같은 조건에서 기존 설계의 적용판과 어떤 성능 차이를 보이는가?}

    \item \textbf{RQ3:}
    \rev{입력 표현, 학습 목적, 단계별 기저위험, 모델 구조는 각각 성능 차이를 얼마나 설명하며, 그 몫은 코호트마다 어떻게 다른가?}
\end{itemize}

\rev{이를 위해 맵, 시점, 팀 규모로 나눈 12개 코호트의 공개 경기 텔레메트리에서 예측 시점 이후의 정보가 섞이지 않는 팀--단계 표본을 대규모로 구성하고, 두 기존 설계를 같은 표본과 평가 조건에 적용하여 비교한 뒤, 설계를 \rev{한 단계씩} 바꾸어 각 단계의 조건부 기여를 평가한다.}

% 연구 질문에 대응하는 기여
본 논문의 기여는 세 질문에 대한 답으로 정리된다.
\begin{itemize}
    \item \textbf{RQ1에 대한 답.}
    \rev{서식지 축소 관점의 해석 가능한 hazard 모델을 제시한다. 이동 부담과 경쟁 밀도가 클수록, 남은 팀원이 적을수록 위험이 높다는 방향은 모든 코호트에서 같고, 구역 안쪽 깊이와 경쟁 밀도의 연관은 단계에 따라 달라진다.}

    \item \textbf{RQ2에 대한 답.}
    \rev{본 설계는 모든 코호트에서 두 적용판보다 AUC가 높고 같은 단계 안의 비교에서도 앞서며, 단계마다 달라지는 위험 수준을 정확히 나타낸다.}

    \item \textbf{RQ3에 대한 답.}
    \rev{\rev{Sagong 적용판과의 차이는 주로 서식지 맥락에서, KRAFTON 적용판과의 차이는 주로 단계별 기저위험과 학습 목적에서 오며,} \rev{모델 구조가 같은 경기 안의 판별을 의미 있게 높인다는 근거는 없었고, 구조는 경기 간 위험 차이에만 작게 기여한다.}}
\end{itemize}
